% !TEX program = xelatex
% ============================================================
% 全国大学生数学建模竞赛(CUMCM) 论文模板 —— 匿名版
% 编译：xelatex main.tex（建议连编两次，生成目录/交叉引用/页码）
% 依赖包（TeX Live / CTeX 完整安装均自带）：
%   ctex geometry setspace amsmath amssymb bm graphicx booktabs
%   caption algorithm algpseudocode listings xcolor fancyhdr hyperref
% 匿名要求：正文不得出现学校、姓名、指导教师、学号、联系方式。
% ============================================================
\documentclass[12pt, a4paper]{ctexart}

% ---------- 页面与行距 ----------
\usepackage[left=2.5cm,right=2.5cm,top=2.5cm,bottom=2.5cm]{geometry}
\usepackage{setspace}
\onehalfspacing

% ---------- 数学 ----------
\usepackage{amsmath,amssymb,bm}

% ---------- 图表 ----------
\usepackage{graphicx}
\usepackage{booktabs}
\usepackage{caption}
\captionsetup[figure]{name=图}
\captionsetup[table]{name=表}
% 约定：表标题在表上方，图标题在图下方（由 \caption 的位置决定）。

% ---------- 算法伪代码 ----------
\usepackage{algorithm}
\usepackage{algpseudocode}
\floatname{algorithm}{算法}

% ---------- 代码附录 ----------
\usepackage{xcolor}
\usepackage{listings}
\lstset{
  basicstyle=\ttfamily\small,
  keywordstyle=\color{blue!70!black},
  commentstyle=\color{gray},
  stringstyle=\color{orange!80!black},
  numbers=left, numberstyle=\tiny\color{gray},
  breaklines=true, frame=single, tabsize=4, showstringspaces=false,
}
% 注意：listings 下的中文注释可能显示异常，建议附录代码注释用英文。

% ---------- 页眉页脚（匿名：不出现任何身份信息）----------
\usepackage{fancyhdr}
\pagestyle{fancy}
\fancyhf{}
\fancyfoot[C]{\thepage}
\renewcommand{\headrulewidth}{0pt}

% ---------- 超链接 ----------
\usepackage[hidelinks]{hyperref}

\title{\textbf{基于 XXX 的 XXX 模型}}   % 点明研究对象+核心方法，勿抄赛题原标题
\author{}   % 匿名：留空，不得写学校/姓名/指导教师
\date{}

\begin{document}

% 国赛提交时的编号页由组委会系统/当年模板生成；正文（本文件）保持匿名。
\maketitle
\thispagestyle{fancy}

% ================= 摘要 =================
\begin{center}\textbf{\zihao{4} 摘\quad 要}\end{center}

% 红线：每个子问题都要有量化结果（每问必有一数）。
% 结构：背景1句 → 总体思路1句 → 各问(方法+数值结果) → 检验/亮点 → 关键词
本文针对……问题，建立了……模型。针对问题一，……，得到……（务必给出具体数值）。
针对问题二，……。针对问题三，……。经灵敏度分析与误差分析验证，模型……。

\vspace{1em}
\noindent\textbf{关键词：}\quad 关键词一；关键词二；关键词三；关键词四

\newpage
\tableofcontents   % 可选，国赛不强制；用了要连编两次
\newpage

% ================= 正文 =================
\section{问题重述}
\subsection{问题背景}
……（用自己的语言概括背景，不要复制题目原文）

\subsection{问题提出}
本文需要解决以下问题：
\begin{itemize}
  \item 问题一：……
  \item 问题二：……
\end{itemize}

\section{问题分析}
% 红线：讲"为什么选这个模型"，而不是"我们将用 A 方法做 X"。
\subsection{问题一分析}
……（数学本质是什么？为什么选这个模型？对比了哪种替代方案？）

\begin{figure}[htbp]
  \centering
  % \includegraphics[width=0.8\textwidth]{figures/flowchart.pdf}
  \caption{整体建模流程图}
  \label{fig:flow}
\end{figure}

\section{模型假设}
% 红线：每条假设都要在后文用到；禁止"数据真实可靠"这类废话假设。
\begin{enumerate}
  \item 假设……。理由：……。
  \item 假设……。理由：……。
\end{enumerate}

\section{符号说明}
\begin{table}[htbp]
  \centering
  \caption{符号说明}
  \begin{tabular}{clc}
    \toprule
    符号 & 含义 & 单位 \\
    \midrule
    $x_i$ & …… & —— \\
    $T$   & …… & s \\
    \bottomrule
  \end{tabular}
  \label{tab:notation}
\end{table}

\section{问题一模型的建立与求解}
% 六步：目标 → 选型理由 → 公式(本题适配) → 参数 → 求解 → 结果(表+图) → 分析
\subsection{模型的建立}
根据……，建立如下模型：
\begin{equation}
  \min\ f(\bm{x})=\sum_{i=1}^{n} c_i x_i, \quad \text{s.t.}\ g(\bm{x})\le 0 .
  \label{eq:model1}
\end{equation}
式中各符号含义见表~\ref{tab:notation}。

\subsection{模型的求解}
采用……算法求解，流程见算法~\ref{alg:1}。
\begin{algorithm}[htbp]
\caption{……求解算法}\label{alg:1}
\begin{algorithmic}[1]
\Require 输入：……
\Ensure 输出：……
\State 初始化……
\While{未收敛}
  \State 更新……
\EndWhile
\State \Return 最优解
\end{algorithmic}
\end{algorithm}

\subsection{结果与分析}
求解得到……（给出具体数值结果）。% 红线：每个图表后必须有解释。

\subsubsection{灵敏度分析}
对关键参数做 $\pm 10\%$、$\pm 20\%$ 扰动，结果表明……，模型稳健。

\section{问题二模型的建立与求解}
（同问题一结构：建立 → 求解 → 结果与分析 → 灵敏度/误差分析。每问末尾嵌检验。）

\section{模型的评价与改进}
% 红线：评价"本模型"，不是评价"所用算法"。
\subsection{模型的优点}
\begin{enumerate}
  \item ……
\end{enumerate}
\subsection{模型的缺点}
\begin{enumerate}
  \item ……（承认本模型具体建模选择的局限，不写放之四海皆准的话）
\end{enumerate}
\subsection{模型的改进方向}
……

% ================= 参考文献 =================
\begin{thebibliography}{99}
% GB/T 7714 格式；红线：≥6 条，禁止 CSDN/知乎/百度百科/AI 工具作为来源。
% 专著[M]：作者. 书名[M]. 版次. 出版地: 出版社, 年.
% 期刊[J]：作者. 题名[J]. 刊名, 年, 卷(期): 起-止页.
\bibitem{ref1} 姜启源, 谢金星, 叶俊. 数学模型[M]. 5版. 北京: 高等教育出版社, 2018.
\bibitem{ref2} 作者甲, 作者乙. 论文题目[J]. 期刊名, 2022, 40(3): 12-20.
\bibitem{ref3} Author C, Author D. Paper title[J]. Journal Name, 2021, 15(2): 100-115.
\end{thebibliography}

% ================= 附录 =================
\newpage
\section*{附录 A\quad 问题一求解源代码（Python）}
\addcontentsline{toc}{section}{附录 A}
% 2026 国赛规则：缺必要源程序或程序不能运行，可能取消评奖资格。放完整可运行代码。
\begin{lstlisting}[language=Python]
import numpy as np

# Solution code for Problem 1.
# Keep code comments in English to avoid CJK glyph issues in listings;
# match variable names to the symbols used in the paper body.
def solve():
    pass
\end{lstlisting}

\end{document}



